% Figure from MATH 493 notes, section 69. Compile with the notes preamble.
\begin{tikzpicture}[>=Stealth, font=\scriptsize, scale=0.95]
\begin{scope}
\clip (-2.6,-2.1) rectangle (2.6,2.1);
\foreach \c in {-3,-2,-1,1,2,3} { \draw[blue!70!black, thin] ({\c-6},-3) -- ({\c+6},3); }
\draw[orange!85!black, very thick] (-6,-3) -- (6,3);
\end{scope}
\draw[->, black!55, thin] (-2.8,0) -- (2.55,0) node[below left] {$x$};
\draw[->, black!55, thin] (0,-2.3) -- (0,2.5) node[above] {$y$};
\node[orange!85!black, anchor=west] at (2.65,1.3) {$U$: $\lambda = 0$};
\foreach \c/\y in {-1/1.8, 1/0.8, 2/0.3, 3/-0.2} { \node[blue!70!black, anchor=west] at (2.65,\y) {$\lambda = \c$}; }
\node[anchor=west, align=left, orange!85!black] at (4.9,1.6) {$U = \mathbb{R}(2, 1)$, a subspace of $V = \mathbb{R}^2$};
\node[anchor=west, align=left, blue!70!black] at (4.9,0.7) {$\lambda(x, y) = x - 2y$ vanishes on $U$:\\its level lines are parallel to $U$};
\node[anchor=west, align=left] at (4.9,-0.4) {$U^\perp = \{\lambda \in V^\vee : \lambda(u) = 0 \text{ for all } u \in U\} = \mathbb{R}\lambda$};
\node[anchor=west, align=left] at (4.9,-1.2) {$\dim U^\perp = \dim V - \dim U = 2 - 1 = 1$};
\end{tikzpicture}
